\section*{Sets and Indices}

\begin{itemize}
    \item $I=\{1,\dots,n\}$: set of key target parts, indexed by $i$. In this instance, $n=4$.
    \item $J=\{0,\dots,n\}$: auxiliary index for the number of destroyed parts in the exact probability calculation.
\end{itemize}

\section*{Parameters}

\begin{itemize}
    \item $d_i$: distance from the airport to part $i$.
    \item $p_i^H$: probability that one heavy bomb destroys part $i$.
    \item $p_i^L$: probability that one light bomb destroys part $i$.
    \item $\bar H$: maximum number of heavy bombs available (\texttt{max\_heavy\_bombs}).
    \item $\bar L$: maximum number of light bombs available (\texttt{max\_light\_bombs}).
    \item $F$: total fuel limit (\texttt{fuel\_limit}).
    \item $k$: minimum number of parts that must be destroyed (\texttt{min\_parts\_destroyed}).
    \item $e^H$: fuel efficiency when carrying a heavy bomb (\texttt{fuel\_efficiency\_heavy\_bomb}).
    \item $e^L$: fuel efficiency when carrying a light bomb (\texttt{fuel\_efficiency\_light\_bomb}).
    \item $e^E$: fuel efficiency when empty (\texttt{fuel\_efficiency\_empty}).
    \item $f^{TL}$: fuel required for takeoff and landing per trip (\texttt{fuel\_takeoff\_landing}).
    \item $f_i^H$: fuel consumed by one heavy-bomb trip assigned to part $i$:
    \[
    f_i^H=\frac{d_i}{e^H}+\frac{d_i}{e^E}+f^{TL}.
    \]
    \item $f_i^L$: fuel consumed by one light-bomb trip assigned to part $i$:
    \[
    f_i^L=\frac{d_i}{e^L}+\frac{d_i}{e^E}+f^{TL}.
    \]
\end{itemize}

For a bomb allocation, define the induced destruction probability of part $i$ as
\[
q_i(x_i^H,x_i^L)=1-(1-p_i^H)^{x_i^H}(1-p_i^L)^{x_i^L}.
\]

\section*{Decision Variables}

\begin{itemize}
    \item $x_i^H$ (code: \texttt{candidate[i][0]} / \texttt{xh}): nonnegative integer number of heavy bombs assigned to part $i$.
    \item $x_i^L$ (code: \texttt{candidate[i][1]} / \texttt{xl}): nonnegative integer number of light bombs assigned to part $i$.
\end{itemize}

Thus,
\[
x_i^H,x_i^L \in \mathbb{Z}_{\ge 0}\qquad \forall i\in I.
\]

\section*{Objective}

Maximize the probability that at least $k$ of the $n$ parts are destroyed:
\[
\max \; \mathbb{P}\!\left(\sum_{i\in I} Y_i \ge k\right),
\]
where the Bernoulli random variables $Y_i$ are assumed independent across parts and satisfy
\[
\mathbb{P}(Y_i=1)=q_i(x_i^H,x_i^L)
=1-(1-p_i^H)^{x_i^H}(1-p_i^L)^{x_i^L}.
\]

Equivalently, the implemented objective is
\[
\max \; \sum_{S\subseteq I:\; |S|\ge k}
\prod_{i\in S} q_i(x_i^H,x_i^L)
\prod_{i\in I\setminus S} \left(1-q_i(x_i^H,x_i^L)\right).
\]

\section*{Constraints}

\begin{align}
\sum_{i\in I} x_i^H &\le \bar H, \\
\sum_{i\in I} x_i^L &\le \bar L, \\
\sum_{i\in I} \left(f_i^H x_i^H + f_i^L x_i^L\right) &\le F, \\
x_i^H,x_i^L &\in \mathbb{Z}_{\ge 0} \qquad \forall i\in I.
\end{align}

\section*{Notes on Implementation}

The code does not formulate or solve a linear, mixed-integer linear, or conic model. Instead, it performs an exact recursive enumeration over all feasible integer allocations $(x_i^H,x_i^L)_{i\in I}$ satisfying the stockpile and fuel limits.

For any enumerated allocation, the part-level destruction probabilities
\[
q_i=1-(1-p_i^H)^{x_i^H}(1-p_i^L)^{x_i^L}
\]
are computed first. Then the probability of destroying at least $k$ parts is evaluated exactly by dynamic programming over the number of successful part destructions. Specifically, if $\mathrm{dp}^{(r)}_j$ denotes the probability of exactly $j$ destroyed parts after processing the first $r$ parts, then
\[
\mathrm{dp}^{(0)}_0=1,\qquad \mathrm{dp}^{(0)}_j=0 \; (j>0),
\]
and for $r=1,\dots,n$,
\[
\mathrm{dp}^{(r)}_j
=
\mathrm{dp}^{(r-1)}_j(1-q_r)
+
\mathbf{1}_{\{j>0\}}\mathrm{dp}^{(r-1)}_{j-1}q_r
\qquad j=0,\dots,r.
\]
The objective value of that allocation is then
\[
\sum_{j=k}^{n} \mathrm{dp}^{(n)}_j.
\]

The recursive search uses remaining-resource bounds at each part to cap enumeration:
\[
x_i^H \le \min\!\left\{\text{remaining heavy bombs},\left\lfloor \frac{\text{remaining fuel}}{f_i^H}\right\rfloor\right\},
\quad
x_i^L \le \min\!\left\{\text{remaining light bombs},\left\lfloor \frac{\text{remaining fuel}}{f_i^L}\right\rfloor\right\},
\]
and accepts any allocation with unused bombs and/or unused fuel. No additional per-part capacity or exactly-$k$-parts-destroyed constraints are imposed.
